\chapter{Component Implementation}

\section{Gateway Transaction Handling}

The implemented gateway centers on controlled transaction state changes. A
merchant payment request becomes a pending payment only after merchant access,
merchant operating status, and QR readiness are validated. A refund request is
accepted only when the original payment is
eligible for a full refund. External payment and refund evidence is stored
before the gateway applies final state transitions so that business outcomes are
always supported by durable evidence.

Four implementation decisions are especially important:

\begin{itemize}
  \item merchant readiness is checked before live transaction creation;
  \item active VietQR receiving account configuration is required before QR
  payment creation;
  \item payment creation uses business-aware duplicate handling so that
  equivalent pending requests can be reused safely;
  \item ambiguous external evidence creates a manual-review path instead of
  silently overriding internal state.
\end{itemize}

\section{VietQR Generation}

VietQR generation is implemented as an application-level payment concern rather
than as a frontend-only display feature. When the Merchant API accepts a
payment, the gateway derives a short QR reference, builds a VietQR transfer
payload from the active merchant QR account, and renders a PNG image as a
base64 data URL. The response therefore contains both machine-readable payload
data and an immediately displayable QR image.

The implementation deliberately applies pilot constraints before generating QR
data. Amounts must be whole-number VND values, the transfer purpose is derived
from the gateway QR reference, and payment creation fails if the merchant has no
active QR receiving account. Duplicate equivalent pending payment requests
reuse the existing stored QR evidence, which keeps checkout retry behavior
predictable for merchant backends.

\section{Demo Merchant Integration}

The local demo merchant is implemented as a small FastAPI package under
\code{backend/demo_merchant}. It serves its own Vietnamese HTML, CSS, and
JavaScript checkout while keeping integration credentials on the server side.
For payment creation it serializes the exact request body, calculates the body
SHA-256 digest, signs the canonical merchant string, and calls the existing
\code{POST /v1/payments} contract without changing the public gateway API.

The bank simulator follows the same separation. A checkout button asks the demo
merchant server to send a timestamped provider-HMAC callback; it does not set
the order status itself. The demo order remains pending with an
\code{AWAITING_WEBHOOK} notification state until the worker sends a merchant
webhook. The receiving route verifies the event id, timestamp window, body
digest, and HMAC, and stores processed event ids so duplicate deliveries return
success without applying the state twice.

The demo merchant uses in-memory state intentionally. This keeps it a clear
reference integration rather than a second business database and allows each
presentation run to reset by restarting the process or calling its demo-only
reset endpoint.

\section{Internal Operations Support}

Internal operations functions are implemented as explicit support workflows
rather than as incidental administrative tools. Internal users can register
merchants, review readiness information, issue or rotate merchant credentials,
configure VietQR receiving accounts, inspect transactions, review audit
history, and resolve operational incidents.

Sensitive internal actions are traceable. When an internal user changes
merchant access, resolves reconciliation issues, or performs recovery actions,
the system preserves the action reason and related evidence so the operational
history remains understandable after the fact.

\section{Merchant Dashboard Support}

Merchant-facing implementation follows the architectural boundary already
defined for separate merchant dashboard access. Human merchant users sign in
with dedicated dashboard accounts, and post-login reads derive merchant scope
from the authenticated session rather than from request parameters.

This design supports the merchant dashboard in four ways:

\begin{itemize}
  \item overview and analytics views can summarize merchant activity for a
  selected reporting window;
  \item detail views can expose payments, refunds, and notification history
  without leaking data across merchants;
  \item profile and credential screens can show useful metadata while keeping
  raw secrets hidden;
  \item password maintenance remains a user-facing self-service workflow rather
  than an operational workaround.
\end{itemize}

\section{Asynchronous Control Functions}

The implementation wires the asynchronous design model into concrete control
functions needed for time-based and retry-based behavior:

\begin{itemize}
  \item expiration handling for overdue pending payments;
  \item creation of merchant notification events after final payment or refund
  outcomes;
  \item notification delivery processing with retry awareness and delivery
  history;
  \item reconciliation support for late, mismatched, or otherwise ambiguous
  external evidence;
  \item manual retry support for eligible failed merchant notifications.
\end{itemize}

These functions are exposed through the worker entry point
\code{python -m app.worker.main}. The worker uses database locking so that
multiple worker processes do not process the same expiration or webhook retry
job at the same time. The functions remain separate from ordinary request
handling so payment outcomes, notification delivery, and operational recovery
stay understandable as related but distinct behaviors.

For the classroom flow the worker interval is set to one second. Its model
registry is loaded before job execution so payment expiration flushes can
resolve every SQLAlchemy foreign-key table, and PostgreSQL advisory locks keep
expiration and due-webhook work safe when more than one worker is running.
